% problem-000357 3 晶体结构与晶胞
\section*{3. 晶体结构与晶胞}
\textit{下列为分问。}

\subsection*{3-1}
该图给出了晶胞中的所有原⼦，除“三原⼦”（中⼼原⼦和与其相连的2个原⼦）外，晶胞的其余原⼦都是 $\mathrm { B _ { 1 2 } }$ 多⾯体中的1个原⼦， $\mathrm { B _ { 1 2 } }$ 多⾯体的其他原⼦都不在晶胞内，均未画出。图中原⼦旁的短棍表⽰该原⼦与其他原⼦相连。若上述“三原⼦”都是碳原⼦，写出碳化硼的化学式。

\subsection*{3-2}
该图有什么旋转轴？有⽆对称中⼼和镜⾯？若有，指出个数及它们在图中的位置（未指出位置不得分）。

（图：）

\subsection*{3-3}
该晶胞的形状属于国际晶体学联合会在1983年定义的布拉维系七种晶胞中的哪⼀种？注：国际晶体学联合会已于 2002 年改称布拉维系(Bravais system)为晶系(lattice system)。

\subsection*{4-1}
分别将 $\mathrm { O _ { 2 } } ,$ 、 $\mathrm { K O } _ { 2 } .$ $\mathrm { B a O _ { 2 } }$ 和 $\mathrm { O _ { 2 } [ A s F _ { 6 } ] }$ 填⼊与O—O键⻓相对应的空格中。

\textit{（原卷此处含表格/仪器试剂表，详见 PDF 或 MD 源文件。）}

\subsection*{4-2}
在配合物A和B中，O_{2}为配体与中⼼⾦属离⼦配位。A 的化学式为 $\mathrm { [ C o _ { 2 } O _ { 2 } ( N H _ { 3 } ) _ { 1 0 } ] ^ { 4 + } }$ ，其O—O的键⻓为 147 pm；B 的化学式为 $\mathrm { C o ( b z a c e n ) P y O _ { 2 } }$ ，其 O—O 的键⻓为 126 pm，Py 是吡啶 $\mathrm { ( C _ { 5 } H _ { 5 } N ) }$ ，bzacen 是四⻮配体 $[ \mathrm { C _ { 6 } H _ { 5 } - C ( O ^ { - } ) = C H - C ( C H _ { 3 } ) = N C H _ { 2 } - } ] _ { 2 }$ 。B具有室温吸氧，加热脱氧的功能，可作为⼈⼯载氧体。画出A和B的结构简图（图中必须明确表明O—O与⾦属离⼦间的空间关系），并分别指出A和B中Co的氧化态。
